\subsection{Skill Catalog}

\subsubsection*{Skill-Families \& Execution Profiles}
\label{subsec:skill-families-profiles}
Dieses Dokument beschreibt die 13 Skill-Familien im CHAL-Skillsystem, jeweils mit:
\begin{itemize}
	\item Konfigurationsachsen (Daten, nicht Code)
	\item Vorschlag für 3 Execution Profiles: \textbf{Melee}, \textbf{Ranged}, \textbf{Caster}
\end{itemize}

Annahme:
\begin{itemize}
	\item \textbf{Melee-Profil} = Vanguard / Bruiser / Piercer
	\item \textbf{Ranged-Profil} = Ranger / Marksman / Saboteur
	\item \textbf{Caster-Profil} = Arcanist / Alchemist / Oracle / Hierophant (+ Summoner mit extra Minion-Tag)
\end{itemize}

\subsubsection{F1 – Spinning Assault}
Namensideen: \emph{Maelstrom-, Cyclone-, Gyro-, Vortex-Modul}

\paragraph{Konfigurationsachsen (Daten)}{}

\begin{itemize}
	\item \texttt{base\_radius} – Grundradius um den Caster.
	\item \texttt{can\_move\_while\_channeling} – true/false.
	\item \texttt{move\_speed\_mod} – Bewegungs-Multiplikator während des Spins.
	\item \texttt{tick\_rate} – Hits pro Sekunde.
	\item \texttt{damage\_per\_tick} – Basis-DMG pro Tick.
	\item \texttt{max\_duration} – maximale Channel-Dauer.
	\item \texttt{startup\_time} – Windup bis der Spin aktiv ist.
	\item \texttt{has\_magnet\_pull} – leichter Sog-Effekt ja/nein.
	\item \texttt{damage\_falloff} – Schaden außen vs. innen (optional).
\end{itemize}

\paragraph{Execution Profiles}{}
\scriptsize %\normalsize
\begin{tabular}{p{3cm} | p{4cm} | p{4cm} | p{4cm}}
	Parameter               & Melee-Profil (Vanguard/Bruiser/Piercer)                    & Ranged-Profil (Ranger/Marksman/Saboteur)                               & Caster-Profil (Arcanist/Alchemist/Oracle/Hierophant)                   \\
	\hline
	\texttt{base\_radius}           & mittel–groß (2.5–3.5)                                      & klein–mittel (2.0–2.5, eher um eine Turret/Zone)                       & mittel (2.5–3.0, magische Orbs um den Caster)                          \\
	\texttt{can\_move\_while\_channeling} & true                                                  & abhängig vom Skill: häufig \textbf{false} (stationäre Turret/Zone)          & true, aber \texttt{move\_speed\_mod} reduziert                                  \\
	\texttt{move\_speed\_mod}        & 80–100 \% (Vanguard langsamer, Bruiser normal, Piercer schneller) & 0–60 \% (stationär oder langsamer Move während Spin-Zone)          & 60–80 \% (Caster kann sich bewegen, aber spürbar langsamer)             \\
	\texttt{tick\_rate}             & 3–5 Hits/s                                                 & 2–3 Hits/s (Turret-/Trap-Feeling)                                      & 2–4 Hits/s (Orbs treffen in Intervallen)                               \\
	\texttt{damage\_per\_tick}       & mittel–hoch, mit Defensive-Boni kombiniert                & mittel, dafür guter Clear-Radius / Utility                             & mittel, skaliert stärker mit Spellpower / Status-Effekten              \\
	\texttt{max\_duration}          & lang (4–6 s), limitiert durch Ressource                    & mittel (3–4 s)                                                         & mittel (3–4 s), ggf. kürzer bei starken Sekundäreffekten               \\
	\texttt{has\_magnet\_pull}       & schwach–mittel (zieht Mobs leicht rein)                    & selten, meistens 0                                                     & optional schwach (Void-/Gravity-Theme)                                 \\
	\hline
\end{tabular}
\normalsize
\subsubsection{F2 – Impact Slam}
Namensideen: \emph{Seismic-, Graviton-, Rupture-, Impact-Module}

\paragraph{Konfigurationsachsen}{}

\begin{itemize}
	\item \texttt{cast\_time} – Windup vor dem Schlag.
	\item \texttt{impact\_radius} – Radius der Primär-Explosion.
	\item \texttt{aftershock\_radius} – Radius zusätzlicher Bodenwelle.
	\item \texttt{aftershock\_delay} – Verzögerung bis Aftershock.
	\item \texttt{leap\_or\_dash\_range} – Distanz, falls Bewegung enthalten ist.
	\item \texttt{damage\_multiplier} – Basisschaden pro Impact.
	\item \texttt{knockback\_force} – Rückstoßstärke.
	\item \texttt{can\_chain\_between\_targets} – springt der Effekt über?
\end{itemize}

\paragraph{Execution Profiles}{}
\scriptsize %\normalsize
\begin{tabular}{p{3cm} | p{4cm} | p{4cm} | p{4cm}}
	Parameter           & Melee-Profil                                             & Ranged-Profil                                              & Caster-Profil                                                    \\
	\hline
	'cast\_range'        & 0-2 (self-weaponrange)                                   & 8-12 (mid-range)                                 & 10-14 (Ground-Target-distnace)  \\
	\texttt{cast\_time}         & kurz–mittel (0.3–0.6 s)                                  & mittel (0.4–0.7 s, Aim + Schuss)                          & mittel–lang (0.5–0.9 s, „Meteor“-Cast)                           \\
	\texttt{impact\_radius}     & groß (3–4)                                               & mittel (2–3)                                               & mittel–groß (3–3.5)                                              \\
	\texttt{aftershock\_radius} & optional, oft etwas kleiner als Impact                   & selten, eher 0–1                                           & häufig aktiv (Earthquake-/Comet-Aftershock)                      \\
	\texttt{aftershock\_delay}  & kurz (0.2–0.4 s)                                        & selten genutzt                                             & klar spürbar (0.4–0.7 s, telegraphed AoE-Zone)                   \\
	\texttt{leap\_or\_dash\_range} & 0–6 (Stampede/EQ-Leap möglich)                           & 0 (Schlag ist am Einschlagpunkt des Projektils)            & 0 (reiner Spell-Impact von oben / aus Distanz)                   \\
	\texttt{damage\_multiplier} & sehr hoch, ggf. Self-Punish (Boneshatter-Trauma)         & hoch, Single-Target fokussiert                             & hoch, oft mehr AoE-Anteil / Elementar-Skalierung                 \\
	\texttt{knockback\_force}   & mittel–hoch (Bruiser/Frontline-Werkzeug)                 & gering (max. kleiner stagger)                             & gering–mittel, eher CC/Slow als harter Knockback                \\
	\hline
\end{tabular}
\normalsize
\subsubsection{F3 – Dash Strike}
Namensideen: \emph{Blink-, Phase-, Vector-Module}

\paragraph{Konfigurationsachsen}{}

\begin{itemize}
	\item \texttt{dash\_range} – Distanz der Bewegung.
	\item \texttt{dash\_duration} – wie schnell die Bewegung ausgeführt wird.
	\item \texttt{hit\_shape} – Linie/Breite + End-Radius.
	\item \texttt{can\_pass\_through\_enemies} – ja/nein.
	\item \texttt{end\_aoe\_radius} – AoE am Landepunkt.
	\item \texttt{max\_charges} – Anzahl möglicher Charges.
	\item \texttt{cooldown} – CD pro Charge.
\end{itemize}

\paragraph{Execution Profiles}{}
\scriptsize %\normalsize
\begin{tabular}{p{3cm} | p{4cm} | p{4cm} | p{4cm}}
	Parameter         & Melee-Profil                                                 & Ranged-Profil                                                  & Caster-Profil                                                \\
	\hline
	\texttt{dash\_range}      & mittel–lang (5–8)                                           & kurz–mittel (4–6, eher Reposition)                            & kurz (3–5, Teleport-/Blink-Range)                            \\
	\texttt{dash\_duration}   & sehr kurz (0.1–0.25 s, snappy)                              & kurz (0.15–0.3 s)                                             & kurz, aber mit Cast-Effekt (0.2–0.35 s)                      \\
	\texttt{hit\_shape}       & breite Linie + kleiner End-Radius                           & schmale Linie, optional Autoshot am Ende                      & kleiner Kreis am Start oder Ende (Arcane Nova / Blink-Nova)  \\
	\texttt{can\_pass\_through\_enemies} & meist true (Gap-closer)                           & teilweise true, aber eher Collision wichtig für Positioning   & true (Blink geht durch Ziele hindurch)                       \\
	\texttt{end\_aoe\_radius}  & mittel (2–2.5)                                              & klein (1–1.5)                                                 & klein–mittel (1.5–2), dafür ggf. starker CC oder Debuff      \\
	\texttt{max\_charges}     & 2–3                                                          & 1–2                                                           & 1–2                                                          \\
	\texttt{cooldown}        & kurz (3–6 s pro Charge)                                     & kurz–mittel (4–8 s)                                           & mittel (6–10 s, weil Mobility + Schaden + Utility kombiniert)\\
	\hline
\end{tabular}
\normalsize
\subsubsection{F4 – Nova Burst}
Namensidee: \emph{Pulse-, Wave-, Radiant-, Bloom-, Nova-Module}

\paragraph{Konfigurationsachsen}{}

\begin{itemize}
	\item \texttt{nova\_radius} – Ausdehnung der Explosion.
	\item \texttt{damage\_falloff} – innen außen Unterschied.
	\item \texttt{cast\_time} – Windup der Nova.
	\item \texttt{can\_be\_targeted\_via\_projectile} – ja/nein.
	\item \texttt{element\_tag} – Fire, Cold, Lightning, Holy, etc.
	\item \texttt{cc\_effect} – Knockback, Slow, Stun-Chance.
	\item \texttt{nova\_pulse\_count} – einmalig oder mehrfach pulsend.
\end{itemize}

\paragraph{Execution Profiles}{}
\scriptsize %\normalsize
\begin{tabular}{p{3cm} | p{4cm} | p{4cm} | p{4cm}}
	Parameter         & Melee-Profil                                             & Ranged-Profil                                                & Caster-Profil                                                    \\
	\hline
	'cast\_range'        & 0-2 (self-weaponrange)                                   & 8-12 (mid-range)                                 & 10-14 (Ground-Target-distnace)  \\
	\texttt{nova\_radius}     & klein–mittel (2–2.5, du stehst mitten drin)             & klein (1.5–2, zentriert auf getroffenem Ziel)               & mittel–groß (2.5–3.5, um den Caster)                             \\
	\texttt{damage\_falloff}  & wenig bis kein Falloff (heftiger Melee-Focus)           & stärkerer Falloff (ST auf Primärziel, AoE für Adds)         & moderater Falloff (gleichmäßiger Spell-DMG im Bereich)          \\
	\texttt{cast\_time}       & kurz (0.2–0.4 s)                                        & sehr kurz (instant bei Treffer / on-hit ausgelöst)          & mittel (0.3–0.6 s)                                              \\
	\texttt{can\_be\_targeted\_via\_projectile} | nein                                     & ja (Nova triggert, wenn Projectile trifft)                  | optional (Nova direkt um Caster ohne Projectile)               \\
	\texttt{cc\_effect}       & leichter Knockback, ggf. Stagger                        & kleiner Slow/Weaken auf getroffenen Feinden                 & starker Slow/Freeze/Chill, ggf. Stun-Chance                    \\
	\texttt{nova\_pulse\_count} | 1 (Burst)                                               & 1 (On-Hit-Burst)                                             & 1–3 Pulse (z. B. „wachsende“ Nova)                             \\
	\hline
\end{tabular}
\normalsize
\subsubsection{F5 – Projectile Barrage}
Namensideen: \emph{Shard-, Lance-, Rail-Module}

\paragraph{Konfigurationsachsen}{}

\begin{itemize}
	\item \texttt{projectile\_count} – 1, 3, 7, …
	\item \texttt{spread\_angle} – 0° (Single) bis 90°+ (Fächer).
	\item \texttt{can\_pierce} – ja/nein.
	\item \texttt{chain\_count} – Kettenanzahl.
	\item \texttt{explode\_on\_hit} – AoE beim Treffer.
	\item \texttt{range} – Flugdistanz.
	\item \texttt{projectile\_speed} – Geschwindigkeit.
	\item \texttt{prioritize\_primary\_target} – Ja/Nein (ST vs. Clear).
	\item \texttt{damage\_tags} – Physical, Lightning, Cold, Chaos, etc.
\end{itemize}

\paragraph{Execution Profiles}{}
\scriptsize %\normalsize
\begin{tabular}{p{3cm} | p{4cm} | p{4cm} | p{4cm}}
	Parameter             & Melee-Profil (Cleave-Linie)                                  & Ranged-Profil (klassischer Schütze)                                   & Caster-Profil (magische Geschosse)                                       \\
	\hline
	\texttt{projectile\_count}    & 3–5 „unsichtbare“ Hit-Segmente in einer Linie                & 1–7 Projektile je nach Skill (Marksman eher 1–3, Ranger 3–7)          & 1–3 magische Geschosse                                                     \\
	\texttt{spread\_angle}        & 45–90° (breiter Cleave nah vor dir)                          & 0–45° (Single bis leichte Streuung)                                    & 0–30° (Fokus auf kontrollierbare Spell-Shots)                             \\
	\texttt{can\_pierce}          & optional (Piercer stark, Vanguard/Bruiser eher selten)       & häufig, v. a. für Clear-Skills                                         & optional, dafür eher Chain/Explode auf Casterseite                        \\
	\texttt{chain\_count}         & 0–1                                                          & 0–2                                                                     & 0–3 (Arc/Lightning-Style)                                                 \\
	\texttt{explode\_on\_hit}      & klein (Mini-AoE entlang der Cleave-Linie)                    & situativ (Explosive Arrow, Kinetic Blast)                              & häufiger, v. a. bei Element-Spells                                        \\
	\texttt{range}               & sehr kurz (2–4, quasi Nahkampf-Reichweite)                   & mittel–lang (8–14)                                                     & mittel–lang (8–12)                                                         \\
	\texttt{projectile\_speed}    & sehr hoch / instant (Cleave simuliert Geschossbahnen)        & hoch (klassische Schüsse)                                              & mittel–hoch (Spells können sichtbar „fliegen“)                            \\
	\texttt{prioritize\_primary\_target} | meist ja (hoher ST-Schaden vorne)                     & Split: Marksman ja, Ranger eher nein (Clear)                           & je nach Spell: ST-Nuke oder kleine, gleichmäßig verteilte Treffer         \\
	\hline
\end{tabular}
\normalsize
\subsubsection{F6 – Heavy Nuke}
Namensidee: \emph{Meteor-, Detoantion-, Warhead-, Burst-Module}

\paragraph{Konfigurationsachsen}{}

\begin{itemize}
	\item \texttt{targeting\_mode} – Self-Buff, Projectile, Ground-Target, Corpse-Target.
	\item \texttt{impact\_radius} – AoE-Größe.
	\item \texttt{delay\_time} – Flug-/Fallzeit.
	\item \texttt{damage\_multiplier} – Basisschaden.
	\item \texttt{splash\_falloff} – Randschaden.
	\item \texttt{self\_cost} – HP-/Ressourcen-Kosten.
	\item \texttt{projectile\_speed} – Geschwindigkeit.
\end{itemize}

\paragraph{Execution Profiles}{}
\scriptsize %\normalsize
\begin{tabular}{p{3cm} | p{4cm} | p{4cm} | p{4cm}}
	Parameter          & Melee-Profil                                               & Ranged-Profil                                             & Caster-Profil                                                       \\
	\hline
	'cast\_range'        & 0-2 (self-weaponrange)                                   & 8-12 (mid-range)                                 & 10-14 (Ground-Target-distnace)  \\
	\texttt{targeting\_mode}   & Self-Buff auf nächsten Schlag oder kurzer Wurf (2–4 Range) & Projectile auf anvisiertes Ziel                           & Ground-Target / AoE-Spell (z. B. „klick auf Boden“)                 \\
	\texttt{impact\_radius}    & mittel (2–3)                                              & klein–mittel (1.5–2.5)                                    & mittel–groß (3–4)                                                   \\
	\texttt{delay\_time}       & sehr kurz (0–0.2 s)                                       & kurz (Flugzeit 0.2–0.4 s)                                 & klar sichtbar (0.4–0.8 s, telegraphed Meteor/Comet)                 \\
	\texttt{damage\_multiplier} | sehr hoch, ggf. mit Self-Punish (Trauma/HP-Kosten)        & hoch, ST mit kleinem Splash                               & hoch–sehr hoch, oft „Glass Cannon“-Spell                           \\
	\texttt{splash\_falloff}   & wenig (nah am Schlag vollen Schaden)                      & moderat                                                   & klarer Falloff (Zentrum sehr stark, Ränder spürbar schwächer)      \\
	\texttt{self\_cost}        & optional hoch (z. B. HP-Opfer für starke Melee-Nuke)      & gering–keine                                              & mittel (Mana-/Ressourcen-Kosten hoch, HP-Kosten selten)            \\
	\texttt{projectile\_speed} | instant/niedrig relevant                                  & hoch                                                      & mittel (Nuke soll „fühlbar“ sein, nicht Hitscan)                    \\
	\hline
\end{tabular}
\normalsize
\subsubsection{F7 – Ground Hazard / DoT-Zone}
Namensideen: \emph{Contagion-, Zone-, Corrosion-, Hazard-, Grid-Module}

\paragraph{Konfigurationsachsen}{}

\begin{itemize}
	\item \texttt{shape} – Circle, Line, Donut, Path.
	\item \texttt{radius\_or\_length} – Größe der Fläche.
	\item \texttt{duration} – Lebensdauer der Zone.
	\item \texttt{tick\_rate} – Hits pro Sekunde.
	\item \texttt{damage\_per\_tick} – Basis-DPS.
	\item \texttt{extra\_effects} – Slow, Chill, Poison, Vuln, etc.
	\item \texttt{placement\_source} – unter Caster, Punktziel, Projectile-Treffer.
	\item \texttt{stacking\_behavior} – können mehrere Zonen stacken?
\end{itemize}

\paragraph{Execution Profiles}{}
\scriptsize %\normalsize
\begin{tabular}{p{3cm} | p{4cm} | p{4cm} | p{4cm}}
	Parameter           & Melee-Profil                                           & Ranged-Profil                                                 & Caster-Profil                                                          \\
	\hline
	\texttt{shape}             & Kreis oder kurzer „Ring“ um den Caster                & Kreis am Einschlagspunkt eines Projektils                     & frei platzierbarer Kreis/Line/Donut                                   \\
	\texttt{radius\_or\_length}  & klein–mittel (2–3)                                    & klein–mittel (2–3)                                            & mittel–groß (3–4)                                                     \\
	\texttt{duration}          & kurz–mittel (2–4 s, aggressiver, hoher DPS)           & mittel (3–5 s)                                                & mittel–lang (4–8 s, mehr Kontrolle, weniger Raw-DPS)                 \\
	\texttt{tick\_rate}         & hoch (3–5/s)                                          & mittel (2–3/s)                                                & 2–4/s je nach Spell                                                   \\
	\texttt{damage\_per\_tick}   & hoch, balanced durch eigene Exposition im Feld        & mittel, Clear-Fokus                                           & mittel–hoch, skaliert stark mit Spellpower/DoT-Boni                  \\
	\texttt{extra\_effects}     & Weak Slow/Armor-Shred/Weaken im Nahbereich            & Slow/Poison/Bleed-Vibes                                       & starke Status (Chill/Freeze/Poison/Vuln/Curse etc.)                  \\
	\texttt{placement\_source}  & entsteht unter/um Caster, oft an Melee-Hit gebunden   & entsteht am Einschlag des ersten/letzten Projektils          & entsteht am gewählten Zielpunkt oder unter Zielgegner                \\
	\texttt{stacking\_behavior} | begrenzt (1–2 Zonen)                                  & 2–3 Zonen möglich                                             & höheres Stack-Potential, aber mit Balance-Grenzen                    \\
	\hline
\end{tabular}
\normalsize
\subsubsection{F8 – Trap / Remote Trigger}
Namensideen: \emph{Mesh-, Ambush-, Trap, Trigger-, Snare-Module}

\paragraph{Konfigurationsachsen}{}

\begin{itemize}
	\item \texttt{trap\_count} – Anzahl gleichzeitig aktiver Fallen.
	\item \texttt{arming\_time} – Zeit bis Falle scharf ist.
	\item \texttt{trigger\_radius} – Auslöseradius.
	\item \texttt{trap\_duration} – Lebensdauer der Falle.
	\item \texttt{damage\_on\_trigger} – Burst-Schaden.
	\item \texttt{remote\_trigger\_option} – darf Spieler manuell auslösen?
	\item \texttt{throw\_range} – wie weit Fallen platziert werden können.
\end{itemize}

\paragraph{Execution Profiles}{}
\scriptsize %\normalsize
\begin{tabular}{p{3cm} | p{4cm} | p{4cm} | p{4cm}}
	Parameter            & Melee-Profil                                            & Ranged-Profil                                              & Caster-Profil                                                    \\
	\hline
	\texttt{trap\_count}         & niedrig–mittel (2–4), du bist eh nah dran              & mittel–hoch (3–6)                                          & mittel (3–5)                                                     \\
	\texttt{arming\_time}        & sehr kurz (0–0.2 s)                                    & kurz (0.2–0.4 s)                                           & kurz–mittel (0.2–0.5 s)                                         \\
	\texttt{trigger\_radius}     & klein–mittel (1.5–2.5)                                 & mittel (2–3)                                               & mittel (2–3)                                                     \\
	\texttt{trap\_duration}      & kurz–mittel (4–8 s, eher offensiv)                     & mittel–lang (6–12 s, Gebietskontrolle)                    & mittel–lang (6–12 s)                                             \\
	\texttt{damage\_on\_trigger}  & sehr hoch, ggf. kleiner Self-Risk, wenn zu nah         & hoch, AoE-Clear                                            & hoch, mit zusätzlichen Status-Effekten                          \\
	\texttt{remote\_trigger\_option} | selten, hauptsächlich Auto-Trigger                  & gelegentlich (Sniper-Detos)                                & häufiger (klassische „Detonate Traps“-Caster)                   \\
	\texttt{throw\_range}        & kurz (2–4, droppt quasi vor deinen Füßen)              & mittel–lang (8–12)                                         & mittel (6–10)                                                    \\
	\hline
\end{tabular}
\normalsize
\subsubsection{F9 – Orbits \& Orbs}
Namensideen: \emph{Orb-, Cluster-, Swarm-, Satellite-Module}

\paragraph{Konfigurationsachsen}{}

\begin{itemize}
	\item \texttt{orb\_count} – Anzahl aktiver Orbs.
	\item \texttt{orbit\_radius} – Abstand zum Zentrum.
	\item \texttt{rotation\_speed} – Umdrehungen pro Sekunde.
	\item \texttt{hit\_mode} – Kontakt-Hit, Auto-Projectiles oder Puls-AoE.
	\item \texttt{duration} – orb-Lebensdauer.
	\item \texttt{detach\_option} – Können Orbs „abgeschossen“ werden?
	\item \texttt{targeting\_logic} – nächster Gegner, zufällig, in Linie etc.
\end{itemize}

\paragraph{Execution Profiles}{}
\scriptsize %\normalsize
\begin{tabular}{p{3cm} | p{4cm} | p{4cm} | p{4cm}}
	Parameter        & Melee-Profil                                              & Ranged-Profil                                                & Caster-Profil                                                   \\
	\hline
	\texttt{orb\_count}      & 1–3 (dicke Orbs nah am Körper)                            & 1–2 pro Turret/Station (mehrere Turrets möglich)            & 3–5 kleinere Orbs                                               \\
	\texttt{orbit\_radius}   & klein–mittel (1.5–2.5, du stehst mit im Hitbereich)       & klein (1–2 um stationäres Objekt)                           & mittel (2–3 um Caster)                                         \\
	\texttt{rotation\_speed} & mittel–hoch (1–2 U/s)                                     & mittel (1 U/s)                                              & variabel (0.8–2 U/s, Spell-Fokus)                              \\
	\texttt{hit\_mode}       & Kontakt-Hits, ggf. leichte AoE-Pulse im Nahbereich       & Auto-Projektile in Zielrichtung / Random innerhalb Range    & Kontakt + Auto-Casts / Lightning-Style                         \\
	\texttt{duration}       & mittel (4–6 s)                                           & lang (8–12 s)                                               & mittel–lang (6–10 s)                                           \\
	\texttt{detach\_option}  & optional (starker Burst, schickt alle Orbs nach vorne)    & selten (Turrets bleiben stationär)                          & häufig optional (Deto-Mechanik für Comet-/Volley-Spells)       \\
	\hline
\end{tabular}
\normalsize
\subsubsection{F10 – Minions \& Swarms}
Namensideen: \emph{Spawn-, Seed-, Hive-, Puppet-, Swarm-Module}

\paragraph{Konfigurationsachsen}{}

\begin{itemize}
	\item \texttt{minion\_count} – Anzahl gleichzeitig beschworener Diener.
	\item \texttt{lifespan} – dauerhaft oder zeitlich begrenzt.
	\item \texttt{ai\_behavior} – melee, ranged, suicide, orbit.
	\item \texttt{summon\_style} – instant, channel, on-kill.
	\item \texttt{control\_mode} – fire\&forget, target-command, follow-only.
	\item \texttt{upkeep\_cost} – Ressourcenkosten / Reservierung.
\end{itemize}

\paragraph{Execution Profiles}{}
\scriptsize %\normalsize
\begin{tabular}{p{3cm} | p{4cm} | p{4cm} | p{4cm}}
	Parameter        & Melee-Profil                                         & Ranged-Profil                                           & Caster-/Summoner-Profil                                                   \\
	\hline
	\texttt{minion\_count}   & klein (1–4, eher „Begleiter“)                        & klein–mittel (1–3 Turrets/Constructs)                   & mittel–hoch (5–12 Skelett-/Geister-Schwärme)                               \\
	\texttt{lifespan}       & lang / dauerhaft                                     & mittel–lang (Turrets halten lange)                      & unterschiedlich: dauerhaft (Skeletons) oder kurz (Raging Spirits/Swarm)   \\
	\texttt{ai\_behavior}    & melee/frontline                                      & ranged/turret                                           & gemischt, aber Fokus auf Spell-/DoT-/Debuff-Patterns                      \\
	\texttt{summon\_style}   & instant, selten Channel                              & instant oder in Salven                                  & instant, Channel oder On-Kill-Procs                                       \\
	\texttt{control\_mode}   & follow + auto-attack                                 & stationär + Auto-Fire                                   & fire\&forget, z. T. Zielsteuerung (Fokus-Target, „angreifen dort“)         \\
	\texttt{upkeep\_cost}    & gering–mittel                                        & mittel                                                  & mittel–hoch (Summoner-Identität, viel Power für Ressourceneinsatz)       \\
	\hline
\end{tabular}
\normalsize
\subsubsection{F11 – Auren \& Shouts}
Namensideen: \emph{Echo-, Emitter-, Induciton-, Aura-Module}

\paragraph{Konfigurationsachsen}{}

\begin{itemize}
	\item \texttt{aura\_radius} – Reichweite.
	\item \texttt{uptime\_mode} – Toggle, Channel, zeitlich begrenzt.
	\item \texttt{buff\_values\_offense} – Schaden/Crit etc.
	\item \texttt{buff\_values\_defense} – Rüstung, Resistenzen, Block etc.
	\item \texttt{pulse\_interval} – Tick-Abstand für Effekte.
	\item \texttt{resource\_reservation} – Reservierter Ressourcenteil.
	\item \texttt{self\_damage\_per\_tick} – für RF-artige Effekte.
\end{itemize}

\paragraph{Execution Profiles}{}
\scriptsize %\normalsize
\begin{tabular}{p{3cm} | p{4cm} | p{4cm} | p{4cm}}
	Parameter              & Melee-Profil                                           & Ranged-Profil                                               & Caster-Profil                                                  \\
	\hline
	\texttt{aura\_radius}          & klein–mittel (4–6), du stehst mitten im Team          & mittel (6–8, mehr „Team-Bubble“ für hintere Reihe)         & mittel–groß (6–10, v. a. für Support-Caster)                   \\
	\texttt{uptime\_mode}          & Shouts (kurz stark) + ggf. 1–2 permanente Auren       & meist Totems/Banner mit mittlerer Dauer                    & viele permanente Auren / Reservierungen                        \\
	\texttt{buff\_values\_offense}  & hoch ST/off-Hand-Boni, Melee-DMG, Leech               & Projectile-DMG, Crit, Accuracy                             & Spell-DMG, DoT, Crit, Elementar / Chaos                        \\
	\texttt{buff\_values\_defense}  & Rüstung, Life, Block, Fortify                         & kleinere Def-Buffs, eher Utility                           & Resistenz-/DR-Buffs, Support-/Heal-Over-Time                   \\
	\texttt{pulse\_interval}       & 0.5–1 s (Warcry/Pulse-Heals etc.)                     & 1–2 s                                                      & 1–2 s                                                          \\
	\texttt{resource\_reservation} | gering–mittel (Frontline braucht aktive Ressourcen)   & gering–mittel                                              & mittel–hoch (klassische „Aura-Reservation“-Mechanik)          \\
	\texttt{self\_damage\_per\_tick} | optional für RF-/Berserker-Auren                      & selten/nicht                                               & bei DoT-/RF-artigen Spells möglich, aber eher Nischenfall      \\
	\hline
\end{tabular}
\normalsize
\subsubsection{F12 – Curses \& Detonations}
Namensideen: \emph{Injection-, Curse-, Hex-, Omen-Module}

\paragraph{Konfigurationsachsen}{}

\begin{itemize}
	\item \texttt{curse\_radius} – AoE-Größe beim Cast.
	\item \texttt{max\_cursed\_targets} – Limit pro Cast/Skill.
	\item \texttt{debuff\_types} – -Res, +DamageTaken, Slow, Weaken, etc.
	\item \texttt{detonate\_condition} – On-Kill, On-Recast, On-Secondary-Skill.
	\item \texttt{detonate\_radius} – AoE beim Deto.
	\item \texttt{scaling\_with\_curse\_stacks} – Skalierungsformel.
	\item \texttt{duration} – Laufzeit des Debuffs.
\end{itemize}

\paragraph{Execution Profiles}{}
\scriptsize %\normalsize
\begin{tabular}{p{3cm} | p{4cm} | p{4cm} | p{4cm}}
	Parameter              & Melee-Profil                                             & Ranged-Profil                                               & Caster/Oracle-Profil                                       \\
	\hline
	\texttt{curse\_radius}         & klein–mittel (2–3, du „markierst“ im Nahkampf)          & klein (1.5–2, Sniper-Marks)                                & mittel–groß (3–4, Flächen-Curse)                           \\
	\texttt{max\_cursed\_targets}   & gering–mittel (3–6)                                     & gering–mittel (3–6)                                        & mittel–hoch (6–12)                                         \\
	\texttt{debuff\_types}         & physischer Vulnerability, Armor-Shred, Bleed-Boost      & Projectile-DMG+/Crit-Marks, -Evasion                       & -Resistenzen, +DoT-Taken, Slow/Weaken                      \\
	\texttt{detonate\_condition}   & On-Kill (Corpse Explosion) oder „Heavy-Hit“-Finisher    & On-Recast / On-Hit mit bestimmtem Skill                    & On-Recast, On-Channel-Ende, spezifische Hexblast-Trigger   \\
	\texttt{detonate\_radius}      & mittel (2–3)                                            & klein–mittel (2–2.5)                                       & mittel–groß (3–4)                                          \\
	\texttt{scaling\_with\_curse\_stacks} | simpel (starker One-Shot, wenig Stack-Logik)      & moderat (Stacks wichtig, aber nicht extrem)                & stark stack-basiert (mehr Stacks → massiv stärkere Detos) \\
	\texttt{duration}             & kurz–mittel (4–8 s)                                     & mittel (6–10 s)                                            & mittel–lang (8–14 s, für Setups)                          \\
	\hline
\end{tabular}
\normalsize
\subsubsection{F13 – Target Marks \& Focus Locks}
Namensideen: \emph{Mark-, Focus-, Lock-Module}

\paragraph{Konfigurationsachsen}{}

\begin{itemize}
	\item \texttt{max\_marked\_targets} – Wie viele Ziele gleichzeitig markiert sein können.
	\item \texttt{mark\_duration} – Dauer der Markierung.
	\item \texttt{apply\_mode} – Wie die Markierung aufgebracht wird (On-Hit, On-Cast, Ground-Target, Blick).
	\item \texttt{consume\_condition} – Wann die Markierung verbraucht wird (nächster Hit, Tod, Deto-Skill, Ally-Hit).
	\item \texttt{effects\_on\_attacker} – Buff-Effekte für Angreifer gegen markierte Ziele (z. B. +Crit, +Damage, On-Hit-Effekte).
	\item \texttt{effects\_on\_victim} – Debuffs auf dem markierten Ziel (z. B. -DamageDealt, +DamageTaken, -MoveSpeed).
	\item \texttt{mark\_visibility} – Wie stark die Markierung visuell hervorgehoben wird (UI/FX-Intensität).
	\item \texttt{recast\_behavior} – Verhalten bei erneutem Wirken (überschreiben, verschieben, refreshen).
\end{itemize}

\paragraph{Execution Profiles}{}
\scriptsize %\normalsize
\begin{tabular}{p{3cm} | p{4cm} | p{4cm} | p{4cm}}
	Parameter             & Melee-Profil (Vanguard/Bruiser/Piercer)                            & Ranged-Profil (Ranger/Marksman/Saboteur)                                & Caster/Oracle-Profil (Arcanist/Alchemist/Oracle/Hierophant)                    \\
	\hline
	\texttt{max\_marked\_targets}  & 1–2 (Fokus auf einzelnes Priority-Target im Nahkampf)              & 1–3 (Boss + 1–2 wichtige Adds)                                           & 1–5 (Oracle/Support kann mehrere Ziele gleichzeitig markieren)                \\
	\texttt{mark\_duration}       & kurz–mittel (4–8 s, aggressiver, hit-naher Playstyle)             & mittel (6–10 s, gut für Pre-Mark vor Engage / Kiten)                     & mittel–lang (8–12 s, als Setup- und Teamwerkzeug)                             \\
	\texttt{apply\_mode}          & primär On-Hit im Nahkampf (z. B. Heavy-Angriff markiert Ziel)      & On-Hit per Projectile oder gezielter Mark-Skill (Single-Target-Aim)      & On-Cast mit freier Zielwahl oder kleiner AoE-Cone/Linie auf Distanz           \\
	\texttt{consume\_condition}   & \texttt{OnNextHitByCaster} (nächster Schlag verbraucht Mark für Big-Hit)  & \texttt{OnHitByCaster} oder \texttt{OnDeath} (Sniper-Focus, On-Kill-Chains)            & \texttt{OnDetonateSkill} oder \texttt{OnHitByAnyAlly} (Team-Synergien / Spell-Detos)        \\
	\texttt{effects\_on\_attacker} | +CritChance/+CritMulti/+Damage vs. Marked, ggf. Life/Res on Hit    & +CritMulti, Penetration, Projectile-Flat-DMG vs. Marked                  & +ResourceGain, +DoTScaling, +StatusChance vs. Marked                          \\
	\texttt{effects\_on\_victim}   & -DamageDealt, -Armor, -MoveSpeed, leichter Stagger vs. Marked      & +DamageTaken vom Mark-Anwender, -Evasion/-Dodge                          & -Resistenzen, +DoTTaken, -DamageDealt, evtl. -Cast-/AttackSpeed               \\
	\texttt{recast\_behavior}     & überschreibt alte Mark, Fokus bleibt meist auf einem Ziel          & Recast verschiebt Mark auf neues Ziel oder überschreibt die älteste Mark | kann Marks refreshen oder mehrere parallel halten (Support-/Control-Rolle)    \\
	\hline
\end{tabular}
\normalsize